% !PRINTED 333
\weberpage{333}
%<<CONT
\[
\begin{aligned}
H' &= (-1)^{n-1}(n-1)\,N(\mathfrak{a})^{-n}(\lambda_1\lambda_2\ldots\lambda_n)^{n-2}, \\
&= (-1)^{n-1}(n-1)\,N(\mathfrak{a})^{-2}\,T^{n-2}
\end{aligned}
\]
and consequently, by (13) and (16),
\weq{17}{H = (-1)^{n-1}(n-1)\,T^{n-2}\Delta.}

\begin{center}
§ 88. \gesperrt{Ideal Classes and Classes of Forms.}
\end{center}

In Bd. II, § 170 we explained the \gesperrt{equivalence} of ideals and the \gesperrt{ideal classes}. According to this, two ideals $\mathfrak{a}$, $\mathfrak{b}$ were equivalent. If there is an integral or fractional \gesperrt{number} of the field $\Omega$ which is the quotient of $\mathfrak{b}$ and $\mathfrak{a}$, that is:
\weq{1}{\eta\,\mathfrak{a} = \mathfrak{b},}
and it has been shown that the ideals are thereby divided into classes, and that the \gesperrt{number of these classes is finite}. In some cases it is expedient to take the concept of equivalence somewhat more narrowly and to call $\mathfrak{a}$ and $\mathfrak{b}$ equivalent only if there is a number $\eta$ with \gesperrt{positive norm} that satisfies condition (1). This difference naturally does not come into consideration when the norms of all numbers are positive, as for example in the imaginary quadratic field. Neither does it come into consideration when there are \gesperrt{units with the norm} $-1$, because, if $\varepsilon$ is such a unit, either $N(\eta)$ or $N(\varepsilon\eta)$ is positive, and $\varepsilon$ and $\varepsilon\eta$ satisfy condition (1) at the same time.

But if there are no units whose norm $=-1$, while there are other numbers in $\Omega$ with negative norm, then under the narrower definition of equivalence every ideal class divides once again into two classes. We wish to hold here to the narrower concept of equivalence, which comes into force, for example, for some \gesperrt{real quadratic fields}\footnote{Namely, when the \gesperrt{Pell} equation $t^2 - D\,u^2 = -4$ has no solution.}.

In (1) $\eta$ may be fractional. But if $\alpha$ is an integer divisible by $\mathfrak{a}$, then $\eta\alpha$ is an \gesperrt{integer divisible by $\mathfrak{b}$}, and from this it follows: If
\weq{2}{(\alpha_1,\ \alpha_2,\ \ldots,\ \alpha_n)}
is a basis of $\mathfrak{a}$, then
\[
(\beta_1,\ \beta_2,\ \ldots,\ \beta_n) = (\alpha_1\eta,\ \alpha_2\eta,\ \ldots,\ \alpha_n\eta)
\]
is a basis of $\mathfrak{b}$, and if the first basis is positive, then (under the narrower concept of equivalence) so is the second.
